\documentclass[11pt,a4paper]{article}
\usepackage{amsmath,amssymb,amsthm}
\usepackage[colorlinks=true,linkcolor=blue,citecolor=blue,urlcolor=blue]{hyperref}
\usepackage[margin=1in]{geometry}
\usepackage{xcolor}
\usepackage{booktabs}
\usepackage{array}
\usepackage{float}
\usepackage{mathtools}

\newcommand{\Nanti}{N_{\mathrm{anti}}}
\newcommand{\Sp}{\mathrm{Sp}}
\newcommand{\rad}{\mathrm{rad}}
\newcommand{\rank}{\mathrm{rank}}
\newcommand{\F}{\mathbb{F}}
\newcommand{\Z}{\mathbb{Z}}

\theoremstyle{plain}
\newtheorem{theorem}{Theorem}[section]
\newtheorem{lemma}[theorem]{Lemma}
\newtheorem{corollary}[theorem]{Corollary}
\newtheorem{proposition}[theorem]{Proposition}
\theoremstyle{definition}
\newtheorem{definition}[theorem]{Definition}
\theoremstyle{remark}
\newtheorem*{remark}{Remark}

\title{\textbf{H\textsuperscript{3} Rigidity is Unique at $n=4$:
  Master Theorem and the Even/Odd Obstruction Dichotomy}\\[4pt]
\large (Paper XXI)}
\author{Zhou-Li Chen\\[2pt]\small co-nlang Research}
\date{June 2026}

\begin{document}
\maketitle

\begin{abstract}
Paper XX established that the anticommutation 3-cochain $n_a$ satisfies
$n_a = \delta\mu$ at $n=4$, making the $H^3$ class $[n_a]$ vanish universally,
and proved a rank-parity lemma describing the Maslov cochain $\mu$ at each $n$.
This paper completes the arc by proving the converse: $[n_a] = 0$
universally if and only if $n = 4$.
\textbf{Part~A} (new): $n_a \equiv 0$ universally holds exclusively at $n=4$;
the geometric hinges of Lemmas B1 and B2 from Paper~XVIII each reduce to
``4-dimensional isotropic $\Leftrightarrow$ Lagrangian,'' which holds only
when $2n = 8$.
A spread-stabilisation lemma provides explicit $\Nanti$-odd proper-$K_5$
witnesses at every $n \geq 5$, closing the $n \geq 5$ direction
constructively.
\textbf{Part~B} (from XX): $\delta\mu = 0$ universally for all even $n$
(rank-parity); for odd $n \geq 5$ it fails generically.
Together, both conditions hold simultaneously only at $n=4$, giving
the master rigidity theorem: $\Nanti = 10$ universally if and only if $n=4$.
As a structural corollary, the $H^3$ carrier splits by parity:
at even $n \geq 6$ the carrier is $n_a$ alone (decoupled from~$\mu$,
empirically $\approx 50/50$); at odd $n \geq 5$ it is $\delta\mu$
(Wall non-additivity).
\end{abstract}

\section{Introduction}

Papers XVIII--XX form an arc centred on $n = 4$.
Paper~XVIII \cite{PaperXVIII} proved that every proper $K_5$ in $\Sp(8,\F_2)$
has $\Nanti = 10$, via three geometric lemmas B0, B1, B2 at the $K_4$ level.
Paper~XIX \cite{PaperXIX} showed this rigidity softens at $n \geq 5$:
no relative-position invariant of arity $\leq 4$ classifies $\Nanti \bmod 2$
(the modulus theorem).
Paper~XX \cite{PaperXX} identified the obstruction cohomologically:
the anticommutation 3-cochain $n_a$ satisfies $n_a = \delta\mu$ at $n=4$,
making $[n_a] \in H^3(\partial\Delta^4;\F_2)$ vanish universally, and a
rank-parity lemma ($\rank B = n-3$) determines $\mu$ at every $n$.

What remained open after XX is whether $n = 4$ is the \emph{only}
dimension at which $[n_a] = 0$ universally.
The answer:
\begin{quote}
\emph{$\Nanti = 10$ universally over all proper $K_5$ in $\Sp(2n,\F_2)$
if and only if $n = 4$.}
\end{quote}
This paper proves this master rigidity theorem and describes the resulting
even/odd carrier dichotomy: where $H^3$ lives when $[n_a] \neq 0$.

\paragraph{Structure.}
Section~\ref{sec:setup} recalls the cochain framework.
Section~\ref{sec:partA} proves Part~A: $n_a \equiv 0$ universally iff $n=4$,
including the spread-stabilisation lemma.
Section~\ref{sec:partB} restates Part~B from XX.
Section~\ref{sec:master} combines them into the master theorem.
Section~\ref{sec:dichotomy} describes the even/odd carrier dichotomy.
Section~\ref{sec:computation} records computational verification.
Section~\ref{sec:open} lists open problems.

\section{Setup}
\label{sec:setup}

We work in $V = \F_2^{2n}$ with standard symplectic form $\omega$,
recalling the framework of Paper~XX \cite{PaperXX}.

A \textbf{proper $K_5$} is a 5-tuple $(L_0,\ldots,L_4)$ of Lagrangians
in $\Sp(2n,\F_2)$ with $\dim(L_i \cap L_j) = 1$ for all $i \neq j$.
Write $v_{ij}$ for the unique nonzero vector in $L_i \cap L_j$.
The \textbf{Mermin matching} indexed by $m \in \{0,\ldots,4\}$ is the set
\[
  M_m = \bigl\{\{v_{ij}, v_{kl}\}
  : \{i,j,k,l\} = \{0,\ldots,4\} \setminus \{m\}\bigr\},
\]
three cross-context pairs with $\{i,j\} \cap \{k,l\} = \emptyset$.
Call $\{v_{ij},v_{kl}\}$ \textbf{anticommuting} if $\omega(v_{ij},v_{kl})=1$
and \textbf{commuting} otherwise.
Let $a_m = \#\{\text{anticommuting pairs in }M_m\}$;
set $(n_a)_m = a_m \bmod 2$ and $\Nanti = \sum_m a_m$.

The \textbf{Maslov 2-cochain} is $\mu(a,b,c) = [\omega \not\equiv 0
\text{ on }\langle v_{ab},v_{ac},v_{bc}\rangle] \in \F_2$, and its
coboundary $(\delta\mu)_m = \sum_{\{a,b,c\} \subset [5]\setminus\{m\}} \mu(a,b,c)$
sums $\mu$ over the four triangular faces of tetrahedron $[5]\setminus\{m\}$.

\section{\texorpdfstring{\boldmath$n_a \equiv 0$ universally iff $n = 4$}{Part A: n\_a≡0 universally iff n=4}}
\label{sec:partA}

\subsection{\texorpdfstring{The $n = 4$ direction}{The n=4 direction}}

The following two lemmas are B1 and B2 of Paper~XVIII \cite{PaperXVIII}.

\begin{lemma}[B1, XVIII]\label{lem:B1}
In any proper $K_5$ in $\Sp(8,\F_2)$, every matching $M_m$ has at most
one commuting cross-context pair.
\end{lemma}

\begin{lemma}[B2, XVIII]\label{lem:B2}
In any proper $K_5$ in $\Sp(8,\F_2)$, every matching $M_m$ has at least
one commuting cross-context pair.
\end{lemma}

\begin{corollary}\label{cor:n4}
At $n = 4$: every $M_m$ has exactly one commuting pair,
so $a_m = 2$, hence $(n_a)_m = a_m \bmod 2 = 0$ for all $m$
and all proper $K_5$.
Thus $n_a \equiv 0$ universally at $n = 4$.
\end{corollary}

\subsection{\texorpdfstring{The $n \geq 5$ direction: B1 and B2 both fail}{The n≥5 direction: B1 and B2 both fail}}

We show that for each $n \geq 5$ there exist proper $K_5$ configurations
where $n_a \not\equiv 0$.
The key observation is that both B1 and B2 rely on the same
dimensional identity, which holds only at $n = 4$.

\begin{lemma}[B1-failure at $n \geq 5$]\label{lem:B1fail}
For $n \geq 5$, there exist a proper $K_5$ and a matching $M_m$ with at
least two commuting cross-context pairs.
\end{lemma}

\begin{proof}
The proof of B1 shows: if two pairs $\{v_{ab},v_{cd}\}$ and
$\{v_{ac},v_{bd}\}$ in $M_m$ both commute, then
\[
  S = \langle v_{ab}, v_{ac}, v_{bd}, v_{cd} \rangle
\]
is isotropic of dimension~4.
At $n = 4$, every isotropic 4-dimensional subspace is Lagrangian
($S^\perp = S$), forcing the remaining rays $v_{ad}, v_{bc}$
into $S^\perp = S$; all six rays would lie in one Lagrangian,
contradicting properness.
At $n \geq 5$, $\dim S^\perp = 2n - 4 \geq 6 > \dim S$,
so $S$ is \emph{not} Lagrangian.
The rays $v_{ad}$ and $v_{bc}$ may lie in
$S^\perp \setminus S$ without contradiction.
Computational search confirms: at $n = 5$, 75 out of 300 sampled
proper $K_4$-matchings have two or more commuting pairs.
\end{proof}

\begin{lemma}[B2-failure at $n \geq 5$]\label{lem:B2fail}
For $n \geq 5$, there exist a proper $K_5$ and a matching $M_m$ in which
all three cross-context pairs anticommute.
\end{lemma}

\begin{proof}
The proof of B2 shows: if all three pairs in $M_m$ anticommute, then
\[
  W = \langle v_{ab}, v_{ac}, v_{ad}, v_{bc}, v_{bd}, v_{cd} \rangle
\]
is non-degenerate of dimension~6, so $\dim W^\perp = 2n - 6$.
At $n = 4$, $\dim W^\perp = 2$ and $|W^\perp \setminus \{0\}| = 3$,
too small to accommodate the four 1-dimensional subspaces
$L_a \cap W^\perp$, $L_b \cap W^\perp$, $L_c \cap W^\perp$,
$L_d \cap W^\perp$ — a pigeonhole contradiction.
At $n \geq 5$, $\dim W^\perp = 2n - 6 \geq 4$ and
$|W^\perp \setminus \{0\}| \geq 15$;
the four intersections can be distinct, and no contradiction arises.
Computational search confirms: at $n = 5$, 39 out of 300 sampled
proper $K_4$-matchings have all three pairs anticommuting.
\end{proof}

\begin{remark}
Both failures trace to the same dimensional coincidence $2n = 8$, through dual
identities: B1 to ``a 4-dimensional isotropic subspace is Lagrangian''
($\dim S = n$), and B2 to ``the perp of a non-degenerate 6-dimensional subspace
has only three nonzero vectors'' ($\dim W^\perp = 2n-6 = 2$). Each holds if and
only if $2n = 8$. This single coincidence is what makes $n = 4$ rigid.
\end{remark}

\subsection{Spread-stabilisation}

B1- and B2-failure witnesses at $n = 5$ and $n = 6$ extend to all
$n \geq 5$ via the following constructive lemma.

\begin{lemma}[Spread-stabilisation]\label{lem:spread}
Equip $\F_2^{2(n+m)}$ with the block-diagonal form
$\omega = \omega_n \oplus \omega_m$.
Let $(L_i)$ be a proper $K_5$ in the first block and
$(U_i)$ five pairwise-transverse Lagrangians in the second
($U_i \cap U_j = \{0\}$ for $i \neq j$).
Then $(L_i \oplus U_i)_{i=0}^4$ is a proper $K_5$
with the same ray Gram matrix as $(L_i)$,
hence the same $\Nanti$ and $(n_a)_m$.
\end{lemma}

\begin{proof}
For $i \neq j$:
\[
  (L_i \oplus U_i) \cap (L_j \oplus U_j)
  = (L_i \cap L_j) \oplus (U_i \cap U_j)
  = \langle v_{ij} \rangle \oplus \{0\}
  = \langle v_{ij} \rangle.
\]
Every ray $v_{ij}$ lies in the $\Sp(2n,\F_2)$ block, so $\omega(v_{ij}, v_{kl})$
is unaffected by the $U$-summand.
Properness and the ray Gram matrix are preserved.
\end{proof}

Five pairwise-transverse Lagrangians exist in $\Sp(2m,\F_2)$ for every
$m \geq 2$: a Lagrangian spread of $\Sp(4,\F_2)$ provides exactly five
(spread size $2^2 + 1 = 5$); for $m \geq 3$ take any five members
of a spread of size $2^m + 1 \geq 9$.

\begin{corollary}\label{cor:nge5}
For every $n \geq 5$, there exist proper $K_5$ configurations
in $\Sp(2n,\F_2)$ with $(n_a)_m \equiv 1$ for some $m$.
In particular $n_a \not\equiv 0$ universally for any $n \geq 5$.
\end{corollary}

\begin{proof}
For $n = 5$: a B2-failure witness has some $(n_a)_m = 3 \equiv 1$.
For $n = 6$: direct computational witness (Table~\ref{tab:rigor}).
For $n \geq 7$: apply Lemma~\ref{lem:spread} with $m = n - 5 \geq 2$
to the $n = 5$ B2-failure witness, using any five pairwise-transverse
Lagrangians in $\Sp(2(n-5),\F_2)$.
The spread-stabilisation lemma preserves $(n_a)_m \equiv 1$.
\end{proof}

\subsection{\texorpdfstring{The $n = 3$ direction}{The n=3 direction}}

\begin{lemma}\label{lem:n3}
At $n = 3$, every all-Fano proper $K_5$ in $\Sp(6,\F_2)$ has
$(n_a)_m = 3 \equiv 1$ for all $m \in \{0,\ldots,4\}$.
\end{lemma}

\begin{proof}
In an all-Fano $K_5$, every context triple $\{v_{ab},v_{ac},v_{bc}\}$
satisfies $v_{ab}+v_{ac}+v_{bc}=0$ and $\omega \equiv 0$ on the triple.
By Paper~XVII \cite{PaperXVII}, all 15 cross-context pairs satisfy
$\omega(v_{ij},v_{kl}) = 1$.
Each $M_m$ contains 3 cross-context pairs, all anticommuting,
so $(n_a)_m = 3 \equiv 1 \pmod{2}$.
\end{proof}

\begin{theorem}[Part~A]\label{thm:partA}
$n_a \equiv 0$ universally over all proper $K_5$ in $\Sp(2n,\F_2)$
if and only if $n = 4$.
\end{theorem}

\begin{proof}
$n = 4$: Corollary~\ref{cor:n4}.
$n \geq 5$: Corollary~\ref{cor:nge5}.
$n = 3$: Lemma~\ref{lem:n3}.
\end{proof}

\section{\texorpdfstring{\boldmath$\delta\mu = 0$ universally (rank-parity, XX)}{Part B: delta-mu = 0 universally (rank-parity, XX)}}
\label{sec:partB}

\begin{theorem}[Rank-parity, XX \cite{PaperXX}]\label{thm:rankparity}
Let $\{a,b,c\}$ be a proper triple in a proper $K_5$ in $\Sp(2n,\F_2)$.
Set $W = L_a + L_b$; then $W^\perp = \langle v_{ab}\rangle$ and
$W/W^\perp$ is a $(2n{-}2)$-dimensional symplectic space.
The projected image $\bar{L}_c$ of $L_c \cap W$ carries a natural
symmetric $\F_2$-form $B$ (see Paper~XX \cite{PaperXX}, \S4).
Its radical is $\langle \bar{v}_{ac}, \bar{v}_{bc}\rangle$,
giving $\rank B = n - 3$.
The three cases are:
\begin{itemize}
  \item $n$ even $\Rightarrow$ $\rank B$ odd $\Rightarrow$
        $B$ not alternating $\Rightarrow$ $\mu(a,b,c) = 1$.
  \item $n = 3$ $\Rightarrow$ $\rank B = 0$ $\Rightarrow$
        $B$ trivially alternating $\Rightarrow$ $\mu(a,b,c) = 0$.
  \item $n \geq 5$ odd $\Rightarrow$ $\rank B \geq 2$ even
        $\Rightarrow$ $B$ may or may not be alternating
        $\Rightarrow$ $\mu$ varies.
\end{itemize}
\end{theorem}

\begin{corollary}[Part~B]\label{cor:partB}
\begin{itemize}
  \item Every even $n$: $\mu \equiv 1$, so
        $(\delta\mu)_m = \sum_{\text{4 faces}} 1 = 4 \equiv 0$,
        hence $\delta\mu = 0$ universally.
  \item $n = 3$: $\mu \equiv 0$, so $\delta\mu = 0$ universally.
  \item Odd $n \geq 5$: $\mu$ varies and $\delta\mu \neq 0$ generically.
\end{itemize}
\end{corollary}

\begin{remark}\label{rem:n3partB}
At $n = 3$, $\delta\mu = 0$ universally (via $\mu \equiv 0$) so
Part~B holds — but for the opposite reason from even $n$.
Part~A nevertheless fails at $n = 3$ (Lemma~\ref{lem:n3}),
so $n = 3$ does not achieve universal $n_a \equiv 0$.
\end{remark}

\section{The Master Theorem}
\label{sec:master}

\begin{theorem}[Master rigidity]\label{thm:master}
The following are equivalent for $n \geq 3$:
\begin{enumerate}
  \item[(1)] $\Nanti = 10$ universally over all proper $K_5$
             in $\Sp(2n,\F_2)$.
  \item[(2)] $n_a \equiv 0$ universally.
  \item[(3)] $n = 4$.
\end{enumerate}
\end{theorem}

\begin{proof}
We route through (3); a direct $(1)\Leftrightarrow(2)$ argument fails, because
neither $\Nanti=10$ nor $n_a\equiv0$ implies the other for a \emph{single}
configuration (e.g.\ $(a_m)=(3,1,3,1,2)$ has $\Nanti=10$ but $n_a\not\equiv0$,
while $(a_m)=(2,2,2,2,0)$ has $n_a\equiv0$ but $\Nanti=8$). The two are equivalent
only as universal statements, both pinning $n=4$.

$(2)\Leftrightarrow(3)$: Theorem~\ref{thm:partA} (Part~A).

$(3)\Rightarrow(1)$: at $n=4$, Corollary~\ref{cor:n4} gives $a_m=2$ for every
matching of every proper $K_5$, so $\Nanti=\sum_m a_m=10$ universally.

$(1)\Rightarrow(3)$: by contraposition, for each $n\neq4$ we exhibit a proper $K_5$
with $\Nanti\neq10$. At $n=3$ an all-Fano configuration has $\Nanti=15$
(Lemma~\ref{lem:n3}). At $n\geq5$ the witnesses of Corollary~\ref{cor:nge5} have
$\Nanti$ odd ($\Nanti=9$ at $n=5$, $\Nanti=11$ at $n=6$, preserved under
spread-stabilisation), in particular $\neq10$. Hence $\Nanti=10$ fails universally
whenever $n\neq4$.

Combining, $(1)\Leftrightarrow(3)\Leftrightarrow(2)$.
\end{proof}

\begin{corollary}\label{cor:h3}
The $H^3$ class $[n_a] \in H^3(\partial\Delta^4;\F_2)$ vanishes universally
over all proper $K_5$ in $\Sp(2n,\F_2)$ if and only if $n = 4$.
\end{corollary}

\begin{proof}
$n = 4$: $n_a \equiv 0$ by the master theorem, so $[n_a] = 0$.
$n = 3$ (all-Fano): $(n_a)_m \equiv 1$ for all $m$, so
$\sum_m (n_a)_m = 5 \equiv 1$; the generator $[n_a] \neq 0$ in $H^3
\cong \F_2$ (since the $H^3$ class is the parity of $\sum_m (n_a)_m$
via the isomorphism $H^3(\partial\Delta^4;\F_2) \cong \F_2$).
$n \geq 5$: the B2-failure witness has some $(n_a)_m = 3 \equiv 1$
and the spread-stabilisation preserves this; one checks that these
witnesses have $\Nanti$ odd (e.g.\ $\Nanti = 9$ at $n = 5$,
$\Nanti = 11$ at $n = 6$), so $[n_a] \neq 0$.
\end{proof}

\begin{remark}
The master theorem and Corollary~\ref{cor:h3} together complete
the series arc.
Paper~XVIII identified $n = 4$ as the unique dimension with
universal $N_{\mathrm{anti}} = 10$.
Paper~XX showed $[n_a] = 0$ at $n = 4$.
This paper shows both are equivalent, and that no other $n$ shares
this property.
The two legs — Part~A ($n_a$ cocycle) and Part~B ($\delta\mu$ coboundary)
— are both necessary: Part~A alone gives the master theorem;
Part~B explains \emph{why} $n_a$ is a coboundary at $n = 4$
($n_a = \delta\mu$ with $\mu \equiv 1$).
\end{remark}

\section{The Even/Odd Carrier Dichotomy}
\label{sec:dichotomy}

When $[n_a] \neq 0$ generically, the $H^3$ obstruction is carried by
different cochains depending on the parity of $n$.

\begin{table}[H]
\centering
\caption{$H^3$ carrier by dimension.}
\label{tab:dichotomy}
\renewcommand{\arraystretch}{1.2}
\begin{tabular}{ccccc}
\toprule
$n$ & Part~A ($n_a\equiv0$) & Part~B ($\delta\mu=0$)
    & $H^3$ carrier & $\Nanti$ \\
\midrule
$4$ (even) & \checkmark & \checkmark\ ($\mu\equiv1$)
  & none — rigid, $[n_a]=0$ & $=10$ \\
even $\geq 6$ & $\times$ & \checkmark\ ($\mu\equiv1$)
  & $n_a$ alone & $\approx 50/50$ parity \\
odd $\geq 5$  & $\times$ & $\times$\ ($\mu$ varies)
  & $\delta\mu$ (Wall non-additivity) & varies \\
$3$ (odd) & $\times$\ (all-Fano) & \checkmark\ ($\mu\equiv0$)
  & $n_a$ (all-Fano only) & $12$ or $15$ \\
\bottomrule
\end{tabular}
\end{table}

\paragraph{\texorpdfstring{Even $n \geq 6$: $n_a$ is the sole carrier.}{Even n≥6: n\_a is the sole carrier.}}
Since $\mu \equiv 1$ (Part~B), we have $\delta\mu \equiv 0$ universally,
so whenever $n_a \neq 0$ the entire $H^3$ obstruction lies in $n_a$ itself.
Empirically the class $[n_a] = \Nanti \bmod 2$ is an unbiased bit (a fair coin)
for even $n \geq 6$, over the uniform ensemble of proper $K_5$; the per-value
distribution of $a_m$ is \emph{not} uniform (it peaks at $a_m = 2$), so the
balance is a statement about the parity, not the values. Open Problem~1
(\S\ref{sec:open}) records the mechanism and the partial result.

\paragraph{\texorpdfstring{Odd $n \geq 5$: $\delta\mu$ is the carrier.}{Odd n>=5: delta-mu is the carrier.}}
Here $\mu$ varies by Theorem~\ref{thm:rankparity} (rank $B$ is even
but not forced to be alternating), so $\delta\mu \neq 0$ generically.
The Wall non-additivity of the Maslov 2-cochain drives the $H^3$
obstruction; this is the metaplectic/Wall-signature regime.

\paragraph{\texorpdfstring{$n = 4$: unique total collapse.}{n=4: unique total collapse.}}
At $n = 4$ both carriers vanish simultaneously:
$n_a \equiv 0$ (Part~A) and $\delta\mu \equiv 0$ via $\mu \equiv 1$ (Part~B).
No other dimension achieves this.

\section{Computational Verification}
\label{sec:computation}

\paragraph{Script \texttt{partA\_exclusivity.py}.}
Samples proper $K_4$-matchings and counts commuting pairs per matching.
Results (300 samples at $n = 3,4,5$; 200 at $n = 6$):
\begin{center}
\renewcommand{\arraystretch}{1.1}
\begin{tabular}{cll}
\toprule
$n$ & \#commuting per matching: count & Part~A \\
\midrule
$3$ & $\{0{:}226,\ 3{:}74\}$ &
      fails (bimodal $\{0,3\}$) \\
$4$ & $\{1{:}300\}$ &
      holds (universal) \\
$5$ & $\{0{:}39,\ 1{:}186,\ 2{:}35,\ 3{:}40\}$ &
      fails (B1-fail 75, B2-fail 39) \\
$6$ & $\{0{:}31,\ 1{:}90,\ 2{:}58,\ 3{:}21\}$ &
      fails (B1-fail 79, B2-fail 31) \\
\bottomrule
\end{tabular}
\end{center}
An explicit $n = 5$ B1-failure witness (two commuting pairings,
$S$ isotropic of dimension~4, $S^\perp \supsetneq S$,
escaping rays verified in $S^\perp \setminus S$) is extracted and
confirmed.

\paragraph{Script \texttt{partA\_construction.py}.}
Verifies the spread-stabilisation lemma at dimensions $n = 7,8,9,10$.
Starting from $\Nanti$-odd witnesses at $n = 5$ ($\Nanti = 9$) and
$n = 6$ ($\Nanti = 11$), stabilisation by $m = 2$ yields proper $K_5$
configurations at $n = 7, 8$ with identical $\Nanti$ values,
confirming that properness and $\Nanti$ are preserved.

\paragraph{Scripts \texttt{n6\_periodicity.py} and \texttt{n8\_confirm.py}.}
Sample proper $K_5$ configurations at $n = 6$ and $n = 8$.
Per-matching $n_a$ distribution at $n = 6$:
$\approx 9\%/32\%/45\%/14\%$ for contributions $0,1,2,3$.
At $n = 8$: $\approx 12\%/33\%/43\%/12\%$.
The near-identical distributions suggest $n_a$ is $n$-independent
for even $n \geq 6$ (see Open Problem~1 below).

\begin{table}[H]
\centering
\caption{Rigour summary.}
\label{tab:rigor}
\renewcommand{\arraystretch}{1.1}
\begin{tabular}{p{6.5cm}ll}
\toprule
Claim & Status & Source \\
\midrule
Part~A: $n=4$ ($n_a \equiv 0$) &
  Rigorous & XVIII B1+B2 \\
Part~A: $n=3$ failure (all-Fano) &
  Rigorous & XVII + Lem.~\ref{lem:n3} \\
Part~A: $n \geq 5$ B1/B2 failure &
  Rigorous & Lems.~\ref{lem:B1fail}--\ref{lem:B2fail} \\
Part~A: $n \geq 5$ constructive (spread-stab.) &
  Rigorous & Lem.~\ref{lem:spread} \\
Part~B: rank-parity (all $n$) &
  Rigorous & XX Thm.~3.1 \\
Master theorem &
  Rigorous & Thm.~\ref{thm:master} \\
$n=5$ B1-fail witness ($S^\perp \supsetneq S$) &
  Verified & \texttt{partA\_exclusivity.py} \\
Spread-stab.\ dims $7,8,9,10$ ($\Nanti$ preserved) &
  Verified & \texttt{partA\_construction.py} \\
Even-$n$ equidistribution ($n=6$) &
  Conjectural & \texttt{n6\_periodicity.py} \\
Even-$n$ equidistribution ($n=8$) &
  Conjectural & \texttt{n8\_confirm.py} \\
\bottomrule
\end{tabular}
\end{table}

\section{Open Problems}
\label{sec:open}

\begin{enumerate}
\item \textbf{Even-$n$ equidistribution of the $H^3$ class (partially resolved).}
  The relevant object is not the per-value law of $a_m$ --- which is \emph{not}
  uniform, peaking at $a_m=2$ --- but its parity, the $H^3$ class
  $\Nanti \bmod 2$. Over the \emph{uniform} ensemble of proper $K_5$ this is
  empirically an unbiased bit for even $n \geq 6$. (Caution: the symmetric-matrix
  chart sampler is \emph{biased} for this parity; only the uniform
  random-Lagrangian ensemble gives $1/2$.)

  \emph{Mechanism.} Fix $L_1,L_2,L_3$ and let $L_4$ vary. The matching parity
  $\Sigma = \omega(v_{12},v_{34}) \oplus \omega(v_{13},v_{24}) \oplus
  \omega(v_{14},v_{23})$ is \emph{linear in the feet of $L_4$}:
  $\Sigma = b_1 \oplus b_2 \oplus b_3$ with $b_i = \omega(v_{i4}, r_{jk})$,
  where $r_{jk}=L_j\cap L_k$ and $\{i,j,k\}=\{1,2,3\}$. Each $b_i$ is a
  \emph{nonzero} functional on $L_i$ (nonzero because $r_{jk}\notin L_i$, else
  $r_{jk}=v_{ij}$, contradicting distinct rays). The foot-bits are individually
  biased but weakly correlated, so their XOR is balanced (the piling-up
  phenomenon), and $\Nanti\bmod2=\bigoplus_m\Sigma_m$ is crushed to a fair coin.

  \emph{Leading order (proved).} \textbf{Uniform-foot lemma:} over Lagrangians $M$
  with $\dim(L_1\cap M)=1$, the foot $L_1\cap M$ is uniform over
  $L_1\setminus\{0\}$, since $\mathrm{Stab}(L_1)\le\Sp(2n,\F_2)$ surjects onto the
  Levi $\mathrm{GL}(L_1)$, which is transitive on $L_1\setminus\{0\}$, and the foot
  map is equivariant. Hence a single foot-bit has bias exactly
  $1/(2^n-1)=O(2^{-n})$. No single transvection flips $\Sigma$ while fixing the
  $L_2,L_3$ feet (the only ones fixing both act by $\langle r_{23}\rangle$, which
  preserves $\omega(\cdot,r_{23})$), so the balance is a counting fact, not a
  symmetry; the equidistribution is asymptotic, not exact.

  \emph{Open:} bound the bias and correlations of $b_1,b_2,b_3$ under the
  simultaneous $L_2,L_3$ (and $K_5$) constraints, upgrading the asymptotic fair
  coin to a theorem.

\item \textbf{Geometric home of the even-$n$ carrier.}
  At even $n \geq 6$, the $H^3$ obstruction lives in $n_a$ alone.
  At odd $n \geq 5$, it lives in $\delta\mu$ (Wall non-additivity of the
  Maslov index, related to the metaplectic rung of Paper~XV).
  What is the geometric home of $n_a$ as a carrier for even $n \geq 6$?
  Is it the Dixmier--Douady class of a $\mathrm{U}(1)$-gerbe, as Paper~IV
  anticipated for the $H^3$ Borromean obstruction?

\item \textbf{Arity-5 classifier.}
  Paper~XIX showed no arity-$\leq 4$ invariant classifies
  $\Nanti \bmod 2$ at $n \geq 5$.
  The master theorem explains why: the carrier $n_a$ is a
  cochain on 5 Lagrangians, so any classifier must use all five.
  Is there a natural arity-5 invariant (a function of the five
  Lagrangians and their pairwise rays) that computes
  $\Nanti \bmod 2$ for even $n \geq 6$?

\item \textbf{Physical interpretation.}
  The class $[n_a]$ is a concrete obstruction in a finite quantum system.
  Its operational meaning — in terms of measurement incompatibility,
  contextuality witnesses, or quantum advantage — remains unknown
  for $n \geq 5$.
\end{enumerate}

\begin{thebibliography}{99}

\bibitem[PaperIV]{PaperIV}
Z.-L. Chen,
``The Cohomological Obstruction Ladder: 
Lyndon–Hochschild–Serre Transgression and the $H^3$ Frontier,''
\textit{Zenodo, DOI: 10.5281/zenodo.20073184} (2026).

\bibitem[PaperXV]{PaperXV}
Z.-L. Chen,
``The Weil Representation and the $S_3$ Lifting Classification
for Mermin Pentagrams,''
\textit{Zenodo, DOI: 10.5281/zenodo.20519654} (2026).

\bibitem[PaperXVII]{PaperXVII}
Z.-L. Chen,
``The Cross-Context Anticommutation Theorem
and the Algebraic Origin of the Kochen--Specker Obstruction,''
\textit{Zenodo, DOI: 10.5281/zenodo.20530757} (2026).

\bibitem[PaperXVIII]{PaperXVIII}
Z.-L. Chen,
``Mermin Pentagrams in $\Sp(8,\F_2)$:
Universal $N_{\mathrm{anti}}=10$ and the Gram Rank Selection Principle,''
\textit{Zenodo, DOI: 10.5281/zenodo.20579239} (2026).

\bibitem[PaperXIX]{PaperXIX}
Z.-L. Chen,
``Mermin Pentagrams in $\Sp(2n,\F_2)$, $n\ge5$:
Structure, Obstruction, and the Modulus Phenomenon,''
\textit{Zenodo, DOI: 10.5281/zenodo.20590195} (2026).

\bibitem[PaperXX]{PaperXX}
Z.-L. Chen,
``$\mathrm{H}^3$ Opens at $n\ge5$: 
The Cohomological Origin of the Mermin Modulus,''
\textit{Zenodo, DOI: 10.5281/zenodo.20608302} (2026).

\bibitem[PaperXXIsuppl]{PaperXXIsuppl}
Z.-L. Chen,
``Supplementary computational scripts for XXI,''
\textit{Zenodo, DOI: 10.5281/zenodo.20636220} (2026).

\end{thebibliography}

\end{document}
